\section{Proof Routes, and What Each Asks of the Author}\label{sec:architecture}

The fourteen sources use few proof methods. We call a method as it
appears in Agda a route; two of the routes are module shapes, copied
from development to development, and the others are patterns. This
section describes them, and for each says what an author who brings a
new theory has to supply and what the machine supplies. The route is
the paper's choice, not ours, and for a new theory it follows what the
author has: an exact synthesis algorithm over the ring gives the
levels route, a candidate normal form the normal-form route, a
decomposition of circuits by wires the induction on the width, and a
complete theory for a nearby language an encoding. Four of the shapes
below are such routes; the other two, products of presentations and
Reidemeister--Schreier ports, compose presentations rather than prove
one from scratch. Section~\ref{sec:example} follows one development
from its gate set to its theorem, and Section~\ref{sec:generated}
says how much was generated and how much was written.

\subsection{Levels and edges}\label{sec:levels}

The two-level developments share one statement of the induction of
Section~\ref{sec:by-levels}, over the algorithm's normal word
\aid{nw}, the action of a word on a matrix, and the level
(Appendix~\ref{sec:levels-details} shows the Agda). \aid{Path w M} says
that $\mathit{nw}(w\cdot M)\,w\approx\mathit{nw}(M)$; \aid{EdgesBelow L}
that every generator is a path from every state when both ends of the
edge are below $L$; \aid{EdgesAt L} the same for the edges out of a
state at $L$ that do not go up; and \aid{EdgeStep}, that
\aid{EdgesBelow L} gives \aid{EdgesAt L} for every $L$, is the Main
Lemma. From it a generic module proves completeness by well-founded
induction on the lexicographic order of levels. The statement covers
every generator, not only the basic ones the papers' Main Lemmas are
stated for, so the papers' one-sentence reduction to the basic
generators is a module of its own; squares are closed by one lemma,
\aid{close}, so many cases of a Main Lemma are one call to it with a
relation and a list of level facts; and the definitions are generic in
the normal word, which is how the proof of the missing case of
Hadamard-$\Pi$ served both its algorithm and Cl\'ement's.

What the author supplies on this route is the choice of ring, the
action of the generators on a column, the algorithm \aid{nw} with the
proof that $\mathit{nw}(M)\cdot M=I$, the level with the proof that a
syllable lowers it, and \aid{EdgeStep}, which is the Main Lemma. The
well-founded induction, the composition of paths, \aid{close} and the
presentation record are the shape, written once and copied. Inside the
Main Lemma the machine does the arithmetic. The routes in these proofs
touch at most eight entries of the pivot column, so a state is a fixed
numerator with a window of entries replaced, a \emph{known state},
whose level is read off the parities of the window; each step of a
route is then one ring identity, $a\pm b = \sqrt2\, c$, about symbolic
entries, which a script generates and a ring solver proves, and where
a case analysis is large it is data, a decision tree checked by
evaluation (Section~\ref{sec:trees}). The relations that a square
needs are derived by the rewriter of Section~\ref{sec:search}, by
search where the paper gives no derivation. The evaluator must never
see the full column with its symbolic entries, because it would try to
compute the level (Section~\ref{sec:engineering}).

\subsection{Normal forms for circuits}\label{sec:nf}

The Clifford and real stabilizer developments follow their papers'
normal forms closely, and share a structure. The normal forms are
inductive types: ladders, $Z$-circuits, $D$-ladders and $X$-circuits,
indexed in the real case by a \emph{wire type} that records whether
$Z$ alone or both $X$ and $Z$ are tracked on a wire. A normal form on
$n+1$ wires is a $Z$-circuit and an $X$-circuit that reach down to the
bottom wire, followed by a normal form on $n$ wires shifted up.
Rewriting a word to normal form pushes one gate at a time through the
parts of a normal form, by a mutual recursion whose tables are the
papers' typed relations; in the Clifford development the tables are
generated from those relations and checked letter by letter, and in
the real case they were written by hand, each entry paired with its
relation so that a wrong entry does not typecheck. Uniqueness is
semantic, by the conjugation action on Pauli operators: in the real
case, a normal form's action on $Z_0$ determines its $Z$-circuit, its
action on $X_0$ the $X$-circuit, and the rest acts one wire up. The
matrix semantics is the layer's circuit end over a commutative ring
with an involution and an element $s$ with $2s^2=1$ that the involution
fixes (Section~\ref{sec:matrix-layer}).

The author supplies the normal forms as inductive types, the push of
one gate through each part of a normal form, and the decoding of a
normal form from its action, which is where the paper's uniqueness
argument goes. This route has no module shape: the author states these
obligations for their own normal form, and the library's lemma that
soundness and a unique normal form give completeness closes the
argument, or a page of direct proof does, as in
Section~\ref{sec:example}. The typed relations that the push uses are
generated from the supplement, and soundness is decided by evaluation.

\subsection{Induction on the width}\label{sec:width}

The CNOT-dihedral and phase-affine developments prove completeness by
induction on the number of wires, with a decomposition that the papers
use in different forms. Every circuit is related to
$D\cdot X_b \cdot L$, a diagonal part, a column of $X$ gates, and a
linear circuit. Two circuits with the same matrix have the same $b$;
their linear parts are equal by a coset tower in the style of
Lafont~\cite{lafont2003boolean}, which needs completeness one wire
down; and their diagonal parts are equal because a diagonal part
splits into a part one wire up and a form on the bottom wire, and that
form is determined by its phase function. The phase-affine
hierarchy~\cite{blake2026phaseaffine} has four levels, and we organized
level~3, which adds $T$, around a functor from circuits of a lower
level to circuits of level~3 one wire up, with the new wire~0 as a
control: the bottom-wire forms of level~3 are $T^t$ followed by images
of a level-1 and a level-2 circuit, whose uniqueness follows from
completeness at levels~1 and~2.

The author supplies the decomposition, the coset tower of the linear
part, and the uniqueness of the bottom-wire form. In the CNOT-dihedral
development the local identities that the tower and the decomposition
use, on two and three wires, were found by breadth-first search and
written out as derivations, soundness is decided on tables, and a
reflective normalizer decides the equations between diagonal
expressions that the decomposition meets. In the phase-affine
development the dimension is symbolic, so soundness is by evaluation on
symbolic labels, and the ring identities of its monomial calculus are
generated by a script.

\subsection{Products, extensions and isomorphisms}\label{sec:products}

Most of the qudit Clifford development is assembled from constructions
of the library~\cite{anonymous2027popl}; the long case analysis is
confined to the symplectic normal form. The semidirect product
presents $N\rtimes H$ from presentations of $N$ and $H$ and a word
$c(h,n)$ over $N$ for each pair of generators, by the relations of
both and $h\,n = c(h,n)\,h$, once the client has shown the action well
defined on the relations. The extension, which need not split, adds to
the relations of a normal subgroup $N$ the conjugates of its
generators by the representatives, and each relation of $G/N$ twisted
by the word in $N$ that its lift picks up; for the qudit Clifford group
the twists are phases in $\Z_p$ computed from the Weyl cocycle, and
only $(SH)^3$ has a nonzero one, $(p^2-1)/8$. And a presentation
theorem transports along an isomorphism of rule sets, translations of
generators to words that carry the relations of each rule set to
equations derivable in the other, with both composites the identity up
to derivation; this is how Figure~1 inherits its theorem from the rule
set where the mathematics was done.

Here the author supplies presentations of the factors, and the
well-definedness of the action or the cocycle on the relations; the
product's presentation and its normal form come from the library. A
variant of a rule set costs the translation of its generators and the
derivations of the rules that differ.

\subsection{Encodings, and the bridge to Theorem 4.4}\label{sec:encoding}

Cl\'ement's real Clifford+$CH$ paper~\cite{clement2026realcliffordch}
proves completeness at each width $n\ge2$ by a back-and-forth argument:
a second language $Y$ with a complete theory, an encoding $e$ from
circuits to $Y$ that preserves the matrix, and a decoding $d$ back
such that $d(e(g)) \approx g$ for every gate (its Lemma~8.7) and $d$
maps every relation of $Y$ to a derivable equation (its Lemma~8.8), so
that two circuits with the same matrix have related encodings, hence
related decodings, hence are related. We proved this argument once,
generically, and instantiated it at $n=2$, where $Y$ is the language
of two-level matrices on four basis vectors, and at $n\ge3$, where $Y$
is the paper's group $P$, whose completeness is the paper's
Theorem~4.10, proved from Theorem~4.4 by the Reidemeister--Schreier
method with four cosets. We formalized that too; its two families of
obligations are exactly the paper's items (a) and (b), and in item (b)
the coset action turned out to be a conjugation of a uniform embedding
by four fixed words, so that each obligation reduces to the image of
one relation.

That left Theorem~4.4: the relations of Figure~6 are complete for
\emph{proper} words, those whose letters have distinct indices, at two
semantics, on 4 basis vectors and on $2^n$. The paper cites
Hadamard-$\Pi$ for it, and our Hadamard-$\Pi$ development proves the
same statement, but the two developments were written independently
and differ in every detail, from the order of the indices in a
generator to the ring and the numbering of the basis. A bridge of
about 600 lines carries one to the other
(Appendix~\ref{sec:bridge-details}): a ring map from $\Z[\sqrt2]$ to
$\Z[\frac1{\sqrt2}]$; a translation $\varphi$ from proper letters to
increasing words, with a map $\psi$ back such that
$\psi(\varphi(u))\approx u$; and a map $\Phi$ from scaled matrices to
matrices over $\Z[\frac1{\sqrt2}]$ that is a monoid homomorphism and
sends the matrix of each generator to the matrix of its translation.
Two proper words with the same matrix then have translations with the
same matrix, which the Hadamard-$\Pi$ development relates, so $\psi$
relates the words. With both instances of Theorem~4.4 proved, the
paper's Theorem~8.9 holds with no hypothesis:
\begin{agdacode}
completeness : ∀ n → Complete n
subpresentationᴼ : ∀ n → (n VRel,_===_) IsSubPresentationOf O (2 ^ n)
\end{agdacode}
Here \aid{O (2 \textasciicircum{} n)} is the orthogonal group
$O_{2^n}(\Z[\frac1{\sqrt2}])$ of the Hadamard-$\Pi$ development. The
development first proved the theorem for a monoid of scaled integer
matrices, which are fast to compute with and have no inverses, and
$\Phi$, which is injective and lands in orthogonal matrices, turns it
into the group statement. So the real Clifford+$CH$ circuits, the real
stabilizer circuits and Figure~6 of Hadamard-$\Pi$ are now all stated
against the same group.

The author supplies the encoding, the decoding, the proof that the
encoding preserves the matrix, the two lemmas of the back-and-forth
argument, and the completeness of the target theory; the argument
itself is one generic module of a hundred lines. Decision procedures on
three qubits pin the transcription of the encoding and of the target
theory's rules before any of this is proved
(Section~\ref{sec:numerical}), and the auxiliary equations of the
decoding, several hundred in Cl\'ement's paper, are proved by steps
that are each justified by completeness one width down, decided by
evaluation (Section~\ref{sec:eval}), or a piece of word algebra.

\subsection{Ports of existing Agda proofs}\label{sec:ports}

Bian and Selinger's Agda proofs for two-qubit Clifford+$T$ and
three-qubit Clifford+$CS$ circuits~\cite{bian2023cliffordt2,bian2023cliffordcs3}
were written against their own tactic library. We rebased that library
on our words and presentations, so that their generated proof terms
port mechanically, and connected their syntactic theorem to the
two-level presentations of Greylyn and of Bian and Selinger; after a
change to the tactic library (Section~\ref{sec:engineering}) and
incremental coset normal forms, the three-qubit development checks
with the default strategy in about five minutes. What a
Reidemeister--Schreier argument asks of the author is the translation
of generators, the coset action, and the two families of obligations,
items (a) and (b) as in Section~\ref{sec:encoding}; the theorem that
these give injectivity is in the library, in the original authors' own
interface.

\subsection{A worked example: real stabilizer circuits}\label{sec:example}

The real stabilizer development~\cite{makary2021real} is the one whose
statement Section~\ref{sec:theorem-form} reads. Here we follow it from
the gate set to the theorem, and say at each step who did the work.
The steps, in order: write the gates and the rules; choose the ring and
the gate matrices, and check the rules against the paper numerically;
prove soundness by evaluation; choose the route and state its
obligations; generate the derivations and tables, or search for them
where the paper gives none; prove the route's mathematics; fill the
record. Its route is the paper's, normal forms.

\paragraph{The gate set and the relations.} The gates are $-1$, $H$, $Z$
and $CZ$, a gate set is a family indexed by arity, and the relations
R1--R16 of the paper's Figure~1 are the constructors of an inductive
family, each stated at the width it is drawn on:
\begin{agdacode}
data Gate : ℕ → Set where
  neg-gate       : Gate 0
  H-gate Z-gate  : Gate 1
  CZ-gate        : Gate 2

data _SRel,_===_ : (n : ℕ) → WRel (Gen n) where
  R₁  : n SRel, neg • neg === ε
  R₄  : (₁₊ n) SRel, (H • Z) ^ 4 === neg
  R₁₃ : (₃₊ n) SRel, CZ • CZ ↑ === CZ ↑ • CZ
  ⋯
\end{agdacode}
Wire 0 is the bottom wire and words are in operator order, which the
module says in its header. \aid{Lift-Relation} adds the structural
rules, and one line per generator names its inverse, so that the
presented monoid is a group. This module is 117 lines, and it is the
first thing a reader compares with the paper.

\paragraph{The matrices.} The semantics is the layer's circuit end over
any commutative ring with an involution and an element $s$ with
$2s^2=1$. The development gives the four gate matrices, $H$ as $s$
times an integer matrix, and reads a generator on a higher wire by the
layer's shift:
\begin{agdacode}
valA : Gen n → Op n
valA (gate₀ neg-gate) = scal (- 1#)
valA (gate₁ H-gate)   = emb (scale s (ιO (ZI.ix hZ)))
valA (gate₁ Z-gate)   = emb (ιO (ZI.ix zZ))
valA (gate₂ CZ-gate)  = emb (ιO (ZI.ix czZ))
valA (g ↥)            = up (valA g)
\end{agdacode}
with \aid{hZ}, \aid{zZ} and \aid{czZ} the integer matrices and
\aid{ιO} their image in the ring. A circuit is the product of its
letters. These 30 lines are the second thing a reader compares with
the paper, and the first thing we checked against it numerically
(Section~\ref{sec:numerical}).

\paragraph{Soundness.} Each relation is checked at the width it is
drawn on by evaluation: the integer matrix of the left side is a power
of 2 times that of the right, and the powers of $s$ account for the
difference. Each of the sixteen checks is \aid{refl}; what was written
is the clause that names the relation's two sides and its power of 2,
127 lines in all, and the structural rules follow from the layer's
lemmas.

\paragraph{The normal form, and the route's obligations.} The paper's
Definitions 4.1--4.10 become inductive types of ladders, $Z$-circuits,
$D$-ladders and $X$-circuits (139 lines). The route then asks for three
things, which the author states for their own normal form, the library
having no module for this route (Section~\ref{sec:nf}):
\begin{agdacode}
normalise-ok : (n : ℕ) (w : Circuit n) → n ⊢ w ≈ ⟦ normalise n w ⟧ⁿ
act-unique   : (N N′ : NF n) → ⟦ N ⟧ⁿ ≗ᵃ ⟦ N′ ⟧ⁿ → unsign N ≡ unsign N′
act-of-mat   : {w v : Circuit n} → ⟦ w ⟧ᴬ ≐ ⟦ v ⟧ᴬ → w ≗ᵃ v
\end{agdacode}
The first is the paper's Proposition~5.7: every circuit equals its
normal form. It is a mutual recursion that pushes one generator at a
time through the three parts of a normal form, 821 lines, and each of
its steps is one typed relation. The second is its Propositions
4.11--4.13: normal forms that act alike on Pauli operators by
conjugation are equal up to their sign, proved by decoding the
$Z$-circuit from the image of $Z_0$ and the $X$-circuit from the image
of $X_0$, 473 lines with the action on Pauli operators. The third
connects the matrix to the action, 467 lines: the matrix of a circuit
conjugates Pauli matrices as the action says, and a Pauli matrix is
determined by its entries. From the three, \aid{completeness} is a
page, and the sign is settled by $-U\neq U$ in a nontrivial ring. These
1{,}900 lines are the core of the paper's mathematics; with the tables,
the axioms as letter lists and the reasoning at a fixed width it comes
to 2{,}541 lines (Table~\ref{tab:example}), and that is what the
author writes.

\paragraph{The typed relations.} The push uses 137 typed relations, one
per gate and kind of normal-form wire, which the paper lists in its
Appendix~B and derives in its supplement, with 19 lemmas. A script,
written for this development and not kept, followed the supplement's
derivations line by line for the lemmas and 90 of the typed relations
and emitted each as a list of rewriter steps (Section~\ref{sec:search});
eight more were derived by normalizing both sides, and two by hand.
With the 137 rules that restate them as the push uses them, these are
264 derivations and 5{,}995 lines, half of the development, checked by
the rewriter (316 lines, copied per development); eleven end in a
corrected relation (Appendix~\ref{sec:misprints}). The tables that
index them by the types of the gates (456 lines) were written by hand,
each entry paired with its relation, so that a wrong entry does not
typecheck.

\paragraph{The theorem.} The last module (130 lines) shows that the
matrix of a circuit is unitary, by evaluation on the gate matrices and
the inverse of a word, and fills the record:
\begin{agdacode}
presentation : {n : ℕ} → (n VRel,_===_) IsSubPresentationOf U (2 ^ n)
presentation {n} = record
  { gl   = grouplike
  ; ⟦_⟧  = ⟦_⟧ᵘ
  ; mono = record { isGroupHomomorphism = ⋯
                  ; injective = λ e → completeness (mat-injective e) } }
\end{agdacode}
The homomorphism is the layer's: a monoid homomorphism between groups
is a group homomorphism, and the matrix of a product is the product of
the matrices. Injectivity is completeness after the layer's
\aid{mat-injective}, which says that the matrix determines the
operator. That line is the last of the four things of
Section~\ref{sec:theorem-form} that the reader checks; the third, the
group, is the layer's and is checked once. The guarantee is the
statement of Section~\ref{sec:theorem-form}: R1--R16 present the
subgroup of $U_{2^n}$ that the four gate matrices generate, over every
ring with the stated properties. That this subgroup is the paper's
real Clifford group, the orthogonal matrices that normalize the Pauli
group, is the generation theorem, proved separately; and two further
modules instantiate the statement at $\D[\omega]$ and at $\D[\sqrt2]$.

\paragraph{The cost.} Table~\ref{tab:example} sorts the 11{,}799 lines
by who wrote them: half are generated, a fifth are the paper's
mathematics, a seventh are results beyond the theorem, and the layer,
the rewriter, the statements and the soundness checks are the rest.
Nearly all of the development's commits fall within four hours, and
the theorem with its $\D[\omega]$ instance typechecks from scratch in
about 75 seconds.

\begin{table}
\caption{The real stabilizer development, 11{,}799 lines, by who wrote
what. ``By hand'' means written by the agent from the paper, not
generated.}
\label{tab:example}
\small
\begin{tabularx}{\linewidth}{@{}Y r Y@{}}
\toprule
Part & Lines & Written by \\
\midrule
The relations (117), the interpretation module (288, of which 30 are the gate matrices), the theorem module (130) & 535 & by hand; the reader checks the relations, the 30 lines of gate matrices and the theorem's line \\
The layer: operators, their laws, the matrices & 661 & once; reused by the Clifford development \\
The rewriter & 316 & one copy per development with derivations \\
264 derivations: the supplement's lemmas and typed relations, and the rules as the push uses them & 5{,}995 & a script, checked by \aid{refl} \\
Soundness of R1--R16 & 127 & \aid{refl}, in a clause naming each relation \\
The paper's mathematics: normal forms, tables, the push, uniqueness, the action on Pauli operators, the bridge from matrices to actions, completeness, the axioms as letter lists, reasoning at a width & 2{,}541 & by hand \\
Further results: Corollary 4.16 with its Pauli calculus, the generation theorem, $\mathbb{F}_7$, the two instantiations & 1{,}624 & by hand \\
\bottomrule
\end{tabularx}
\end{table}

\subsection{What is generated and what is written}\label{sec:generated}

Table~\ref{tab:generated} in Appendix~\ref{sec:appendix-routes} gives,
for each development, how many lines were generated by a script and
what they are: 5{,}995 of 11{,}799 for the real stabilizer development,
7{,}347 of 11{,}349 for Clifford, 4{,}337 of 15{,}528 for
$O_n(\Z[\frac12])$ and $L_n$, 2{,}203 of 10{,}406 for Cl\'ement's thesis,
and none for the qudit Clifford and minimality developments. What is
generated is what the paper itself presents as a computation; what is
written is what it presents as an idea. The share is high where the
paper's derivations were followed line by line, and low where its
arguments were replaced by structure, as in the real Clifford+$CH$
development, whose bulk is word algebra written by the agent. Nothing
generated is trusted: the script can be wrong, and Agda checks its
output as it checks everything else.
